
\pagestyle{empty}\thispagestyle{empty}
\vspace*{\fill}
\begin{center}
\fontsize{52}{52}\selectfont
\textsc{scratch au collège}

  \hfil
  
  \includegraphics[width = 5cm]{logoScratchAuCollege3}
\end{center}
\vfill
\begin{center}
\huge
\textsc{algorithmes \  et \  programmation}
\end{center}
\begin{center}
\LogoExoSept{2}
\end{center}
\clearemptydoublepage
%\clearpage

\thispagestyle{empty}

\vspace*{\fill}
\vspace*{-5ex}

\section*{À la découverte du code}

\vspace*{-2ex}

Scratch est un logiciel idéal pour apprendre à programmer. Il a été spécialement conçu pour les enfants et les débutants. La programmation avec Scratch est ludique car il est facile de faire de beaux dessins et des petits jeux. En plus, la programmation est simple, car il suffit de déplacer des blocs pour écrire son code.

\smallskip 

Pourquoi apprendre à coder ? Pour utiliser un ordinateur, je n'ai pas besoin de savoir le programmer ! C'est comme pour les voitures, je n'ai pas besoin de connaître la mécanique pour conduire. Mais, dans le monde qui nous entoure, l'informatique est partout, dans les ordinateurs bien sûr, mais aussi dans nos téléphones, en fait dans tous les appareils électroniques et bientôt ce seront les ordinateurs qui piloteront les voitures ! Il est donc indispensable d'apprendre à parler le langage des ordinateurs. 

\smallskip

Les langages pour programmer un ordinateur sont nombreux, mais une fois qu'un langage est bien compris, les autres s'apprennent plus vite.
Scratch est facile à prendre en main et il permet d'aborder bon nombre de situations de programmation.
Avec Scratch, la programmation devient un jeu et votre ordinateur un compagnon.

\vspace*{-2ex}

\section*{À la découverte des algorithmes}

\vspace*{-2ex}

Un algorithme est une suite d'instructions données permettant d'atteindre un objectif ou de résoudre un problème, un peu comme une recette de cuisine. Comment effectuer une multiplication ? Comment trier une liste ? Quel est le plus court chemin entre deux villes ?

\smallskip

Un algorithme n'est pas lié à un langage, ni même aux ordinateurs ! C'est pourquoi on peut très bien comprendre un algorithme en travaillant sur feuilles. Travailler sur feuilles pour faire de l'informatique, l'idée est surprenante. Mais ce travail permet d'abord de préparer ou de consolider les connaissances apprises devant la machine. Il permet également d'étudier des concepts difficiles à programmer comme par exemple des algorithmes graphiques ou bien encore qui portent sur les mots. 

\smallskip

Alan Turing est un personnage emblématique qui a été l'un des premiers à faire le lien entre travail théorique et travail sur machine. D'une part, il a participé activement à la création d'un des premiers ordinateurs permettant ainsi de décrypter des messages secrets durant la seconde guerre mondiale. D'autre part, il a conçu \emph{la machine de Turing}, encore utilisée de nos jours, qui n'est pas une véritable machine, mais un modèle d'ordinateur sur papier ! 

\medskip

\hfill Alors, prêt à programmer ?

\bigskip
\vspace*{\fill}


\begin{minipage}{0.97\linewidth}
\begin{itemize}
  \item Des vidéos d'explications vous guident sur la chaîne \emph{Youtube} :

  \centerline{\href{https://www.youtube.com/c/ScratchAuCollege}{youtube.com/ScratchAuCollege}}
  
  \item Les scripts Scratch solutions et le livre en couleurs sont disponibles :

  \centerline{\href{https://exo7math.github.io/scratch3-exo7/}{exo7math.github.io/scratch3-exo7}}

  \item  Pour progresser plus vite vous pouvez alterner un chapitre  d'activités Scratch (partie 1), des énigmes sur ce chapitre (partie 2), des activités sur feuilles (elles sont indépendantes les unes des autres, partie 3) et encore des énigmes (partie 4).
\end{itemize}
\end{minipage}

%\vspace*{\fill}

%\newpage
\cleardoublepage
\thispagestyle{empty}
\addtocontents{toc}{\protect\setcounter{tocdepth}{1}}
\tableofcontents
